% =====================================================================
\chapter{Variables aleatorias discretas}
\label{tema:2}
% =====================================================================

Una variable aleatoria asigna un número a cada resultado de un experimento. Este tema estudia las variables discretas, entre ellas la distribución categórica que produce un LLM en cada paso, y sus dos resúmenes principales, la esperanza y la varianza. Se sigue el capítulo~2 y la §4.2 de~\cite{bt}, las lecciones~5 a~7 de~\cite{mit6041} y las clases~4, 5a y~7 de~\cite{mit1805}; la orientación al aprendizaje automático procede de~\cite[§6.2 y §6.4]{mml}.

\section{Variables aleatorias y funciones de masa}

\begin{definicion}[Variable aleatoria discreta]\label{def:va}
Una \emph{variable aleatoria} es una función $X\colon\Omega\to\R$ tal que $\{\omega: X(\omega)\le x\}\in\mathcal{F}$ para todo $x\in\R$. Es \emph{discreta} si toma un conjunto finito o numerable de valores. Su \emph{función de masa} es
\[
p_X(x)=\PP(X=x),
\]
que cumple $p_X(x)\ge 0$ y $\sum_x p_X(x)=1$. Si $X$ y $Y$ tienen la misma función de masa se dice que tienen la misma \emph{distribución}.
\end{definicion}

\begin{definicion}[Distribuciones discretas básicas]\label{def:discretas}
\begin{enumerate}
  \item $X\sim\Ber(p)$, con $p\in[0,1]$, si $X\in\{0,1\}$ y $p_X(x)=p^{x}(1-p)^{1-x}$. Modela un experimento con dos resultados: acierto o fallo.
  \item $X\sim\Bin(n,p)$ si $X$ es el número de éxitos en $n$ ensayos de Bernoulli independientes con probabilidad $p$:
  \[
  p_X(k)=\binom{n}{k}p^{k}(1-p)^{n-k},\qquad k=0,1,\dots,n.
  \]
  \item $X\sim\Cat(p_1,\dots,p_K)$, con $p_k\ge 0$ y $\sum_k p_k=1$, si $X\in\{1,\dots,K\}$ y $\PP(X=k)=p_k$. Equivalentemente, $p_X(k)=\prod_{j=1}^{K}p_j^{\mathbf{1}[j=k]}$.
\end{enumerate}
\end{definicion}

La categórica es la generalización de la Bernoulli a $K$ resultados, y es exactamente lo que produce un LLM en cada posición: la softmax convierte los $K$ logits en el vector $(p_1,\dots,p_K)$ de una categórica sobre el vocabulario. Los valores $1,\dots,K$ son etiquetas sin orden ni distancia: el token~517 no está «entre» el~516 y el~518.

\addtocounter{definicion}{1}%

\section{Esperanza}

\begin{definicion}[Esperanza]\label{def:esperanza}
La \emph{esperanza} o \emph{media} de una variable discreta $X$ es
\[
\E[X]=\sum_x x\,p_X(x),
\]
siempre que la serie converja absolutamente.
\end{definicion}

La esperanza es el centro de gravedad de la función de masa. Para una categórica con etiquetas arbitrarias no tiene sentido «el token medio», pero sí la esperanza de cualquier función numérica de los tokens, por ejemplo de la pérdida.

\begin{teorema}[Esperanza de una función]\label{teo:lotus}
Si $g\colon\R\to\R$ y $X$ es discreta, entonces $\E[g(X)]=\sum_x g(x)\,p_X(x)$.
\end{teorema}

\begin{proposicion}[Linealidad]\label{prop:linealidad}
Para variables $X$, $Y$ con esperanza finita y constantes $a,b,c$,
\[
\E[aX+bY+c]=a\,\E[X]+b\,\E[Y]+c.
\]
La igualdad no requiere independencia.
\end{proposicion}

\addtocounter{definicion}{1}%

\section{Varianza}

\begin{definicion}[Varianza]\label{def:varianza}
La \emph{varianza} de $X$ es $\Var(X)=\E\bigl[(X-\E[X])^{2}\bigr]$, y su \emph{desviación típica} es $\sigma_X=\sqrt{\Var(X)}$, que tiene las mismas unidades que $X$.
\end{definicion}

\begin{proposicion}\label{prop:varianza}
\begin{enumerate}
  \item $\Var(X)=\E[X^{2}]-\E[X]^{2}$.
  \item $\Var(aX+b)=a^{2}\Var(X)$ para constantes $a,b$.
  \item $\Var(X)\ge 0$, con igualdad si y solo si $X$ es constante con probabilidad~1.
\end{enumerate}
\end{proposicion}

\addtocounter{definicion}{1}%

Para $X\sim\Ber(p)$, $X^{2}=X$, luego $\E[X]=\E[X^{2}]=p$ y $\Var(X)=p-p^{2}=p(1-p)$, máxima en $p=1/2$.

\section{Varias variables aleatorias}

\begin{definicion}[Distribución conjunta e independencia]\label{def:conjunta}
La \emph{función de masa conjunta} de $X$ e $Y$ es $p_{X,Y}(x,y)=\PP(X=x,\,Y=y)$; las \emph{marginales} se obtienen sumando, $p_X(x)=\sum_y p_{X,Y}(x,y)$. Las variables $X$ e $Y$ son \emph{independientes} si $p_{X,Y}(x,y)=p_X(x)\,p_Y(y)$ para todo $(x,y)$, y $X_1,\dots,X_n$ son independientes si la conjunta factoriza en el producto de las $n$ marginales.
\end{definicion}

\begin{definicion}[Covarianza]
La \emph{covarianza} de $X$ e $Y$ es $\Cov(X,Y)=\E\bigl[(X-\E[X])(Y-\E[Y])\bigr]=\E[XY]-\E[X]\,\E[Y]$. Si $\Cov(X,Y)=0$, las variables se llaman \emph{incorreladas}.
\end{definicion}

\begin{proposicion}\label{prop:suma}
\begin{enumerate}
  \item Si $X$ e $Y$ son independientes, $\E[XY]=\E[X]\,\E[Y]$ y, por tanto, $\Cov(X,Y)=0$.
  \item $\Var(X+Y)=\Var(X)+\Var(Y)+2\Cov(X,Y)$. En particular, la varianza de una suma de variables independientes es la suma de las varianzas.
\end{enumerate}
\end{proposicion}

El recíproco de (i) es falso: hay variables incorreladas que no son independientes (\cref{eje:2-covarianza}).

\begin{corolario}[Varianza de la media muestral]\label{cor:media}
Sean $X_1,\dots,X_n$ independientes con media $\mu$ y varianza $\sigma^{2}$, y sea $\Xbar_n=\frac1n\sum_{i=1}^{n}X_i$. Entonces $\E[\Xbar_n]=\mu$ y $\Var(\Xbar_n)=\sigma^{2}/n$.
\end{corolario}

El \cref{cor:media} es el resultado central del módulo: promediar $n$ observaciones independientes divide la varianza por $n$ y la desviación típica por $\sqrt n$. El tema~5 lo convierte en garantías sobre el error de cualquier promedio. Como aplicación inmediata, si $X\sim\Bin(n,p)$ se escribe $X=\sum_{i=1}^{n}B_i$ con $B_i\sim\Ber(p)$ independientes, y se obtiene $\E[X]=np$ y $\Var(X)=np(1-p)$.

\begin{table}[h]
\centering\small
\caption{Distribuciones de los temas~2 y~3 y uso en un LLM.}\label{tab:distribuciones}
\begin{tabular}{@{}lllll@{}}
\toprule
Distribución & Valores & $\E[X]$ & $\Var(X)$ & Uso en un LLM \\
\midrule
$\Ber(p)$ & $\{0,1\}$ & $p$ & $p(1-p)$ & acierto en una pregunta de evaluación \\
$\Bin(n,p)$ & $\{0,\dots,n\}$ & $np$ & $np(1-p)$ & aciertos en un banco de pruebas \\
$\Cat(p_1,\dots,p_K)$ & $\{1,\dots,K\}$ & --- & --- & token siguiente \\
$\Unif(a,b)$ & $[a,b]$ & $(a+b)/2$ & $(b-a)^{2}/12$ & número aleatorio del muestreo \\
$\Expo(\lambda)$ & $[0,\infty)$ & $1/\lambda$ & $1/\lambda^{2}$ & tiempos de espera \\
$\N(\mu,\sigma^{2})$ & $\R$ & $\mu$ & $\sigma^{2}$ & pesos iniciales, ruido \\
\bottomrule
\end{tabular}
\end{table}

\section{Muestreo de una distribución categórica}

Para generar un token, el LLM necesita extraer una muestra de $\Cat(p_1,\dots,p_K)$ a partir de un número aleatorio uniforme. Sea $F_k=p_1+\dots+p_k$ la probabilidad acumulada, con $F_0=0$ y $F_K=1$.

\begin{proposicion}[Transformada inversa discreta]\label{prop:inversa-discreta}
Sea $U\sim\Unif(0,1)$, es decir, $\PP(a<U\le b)=b-a$ para $0\le a\le b\le 1$ (\cref{tema:3}). La variable
\[
X=\min\{k: U\le F_k\}
\]
tiene distribución $\Cat(p_1,\dots,p_K)$.
\end{proposicion}

\addtocounter{figure}{1}%

\addtocounter{definicion}{1}%

\addtocounter{definicion}{1}%

